\section{最终推荐特征}

\subsection{推荐表来源}

完整机器可读推荐表位于：

\begin{center}
\path{../summary/recommended_features.csv}
\end{center}

该表包含 dataset、analysis scope、task、feature、feature family、recommendation level、rank/score/plot evidence、label-source 标记、caveat 和 final decision。

LaTeX 中的 compact 表只摘取核心结论。

真正做 baseline planning 时，应以 CSV 为准。

\begingroup
\scriptsize
\setlength{\tabcolsep}{3pt}
\renewcommand{\arraystretch}{1.18}
\begin{longtable}{@{}p{1.7cm}p{1.9cm}p{4.2cm}p{2.3cm}p{4.8cm}@{}}
\caption{两个数据集上的最终推荐特征族}
\label{tab:final_recommended_features}\\
\toprule
数据集 & 任务 & A 级特征 & Reference & 说明 \\
\midrule
\endfirsthead
\toprule
数据集 & 任务 & A 级特征 & Reference & 说明 \\
\midrule
\endhead
XJTU-SY & RUL & \feature{mag__time__mean}, \feature{mag__time__mean_abs}, \feature{mag__time__std} & \feature{mag__time__rms} & 幅值和能量类时域特征在主划分、工况内和跨工况检查中都较稳定。\\
XJTU-SY & Health State & \feature{mag__time__mean}, \feature{mag__time__mean_abs}, \feature{mag__time__std} & \feature{mag__time__rms} & 主划分中水平通道较强，但跨工况检查中 magnitude 特征更稳。\\
XJTU-SY & Early Fault & \feature{mag__time__mean}, \feature{mag__time__mean_abs} & \feature{mag__time__rms} & 早期故障最工况敏感，horizontal 和 spectral 特征只作条件性辅助。\\
PHM2012 & RUL & \feature{h__time__mean_abs}, \feature{mag__time__mean}, \feature{mag__time__mean_abs} & \feature{mag__time__rms} & 水平通道与 magnitude 幅值特征主导。\\
PHM2012 & Health State & \feature{h__time__mean_abs}, \feature{h__time__std}, \feature{h__time__rms} & \feature{mag__time__rms} & 水平通道幅值类特征分离健康状态最强。\\
PHM2012 & Early Fault & \feature{h__time__mean_abs}, \feature{mag__time__mean}, \feature{mag__time__mean_abs} & \feature{mag__time__rms} & 幅值类特征主导，频域特征只作为 secondary 信号。\\
\bottomrule
\end{longtable}
\endgroup


\subsection{A/B/C/Reference 的使用边界}

A 级特征是主线候选。

它们通常同时满足：主划分排名靠前、稳定性检查支持、不是 label-source、物理解释清晰。

B 级特征是 secondary 候选。

它们可能与 A 级高度冗余，或者在某些 split 上略弱，但仍然能为模型提供稳定信号。

C 级特征是条件性候选。

它们适合做 ablation、诊断或工况敏感分析，不适合直接作为跨数据集主线。

Reference 特征主要是 \feature{mag__time__rms}。

它用于检查 HI/FPT 标签构造是否和振动能量一致。

它不用于证明 Health State 或 Early Fault 的独立可检测性。

\subsection{RUL 推荐说明}

RUL 的跨数据集共同推荐是 amplitude / energy-like features。

XJTU-SY 的 A 级是：

\begin{itemize}[leftmargin=2em]
  \item \feature{mag__time__mean}
  \item \feature{mag__time__mean_abs}
  \item \feature{mag__time__std}
\end{itemize}

PHM2012 的 A 级是：

\begin{itemize}[leftmargin=2em]
  \item \feature{h__time__mean_abs}
  \item \feature{mag__time__mean}
  \item \feature{mag__time__mean_abs}
\end{itemize}

这些特征的依据包括 RUL correlation、degradation score、curves 和稳定性检查。

RUL compact subset 可以选择非 label-source 的 magnitude/horizontal amplitude 特征。

\feature{mag__time__rms} 可以作为 reference ablation，但不应成为唯一主特征。

\subsection{Health State 推荐说明}

Health State 的共同模式也是 amplitude-driven。

XJTU-SY 最终更推荐 magnitude features，因为 cross-condition 后它们更稳。

PHM2012 更推荐 horizontal features，因为 official split 中 \feature{h__time__mean_abs}、\feature{h__time__std}、\feature{h__time__rms} 的 boxplot 和统计指标都很强。

Health State 的 baseline 设计应同时考虑两个版本：

\begin{itemize}[leftmargin=2em]
  \item dataset-specific compact set：XJTU 用 magnitude，PHM 用 horizontal；
  \item unified compact set：magnitude + horizontal amplitude 共同输入，让模型自己学习权重。
\end{itemize}

如果包含 \feature{mag__time__rms}，需要在报告里标明它是 Reference。

\subsection{Early Fault 推荐说明}

Early Fault 不能只看主划分 top rank。

PHM2012 中 amplitude features 较稳定。

XJTU-SY 中 Early Fault 工况敏感。

因此推荐策略是：

\begin{itemize}[leftmargin=2em]
  \item PHM2012：以 \feature{h__time__mean_abs}、\feature{mag__time__mean}、\feature{mag__time__mean_abs} 为 A 级；
  \item XJTU-SY：以 \feature{mag__time__mean}、\feature{mag__time__mean_abs} 为 A 级；
  \item XJTU-SY：把 horizontal amplitude、spectral entropy、peak-to-peak/shock-like 特征作为 C 级 ablation；
  \item 两个数据集：\feature{mag__time__rms} 只作 Reference。
\end{itemize}

Early Fault baseline 的报告必须包含 caveat。

特别是 XJTU-SY，如果某个模型在主 split 上表现好，还需要 condition-wise 或 cross-condition 结果支撑。

\subsection{Compact Baseline Subset}

\begin{table}[htbp]
\centering
\caption{下游 baseline 的 compact feature subset 建议}
\label{tab:compact_baseline_suggestions}
\scriptsize
\begin{tabularx}{\textwidth}{lYY}
\toprule
用途 & 建议特征 & 使用边界 \\
\midrule
Full manual set & \config{manual_basic} 全部 45 个特征 & 作为默认 baseline 输入，需 train-only scaling。\\
Compact RUL & \feature{mag__time__mean}, \feature{mag__time__mean_abs}, \feature{mag__time__std}, \feature{h__time__mean_abs}, \feature{h__time__std}, \feature{v__time__std} & 排除或单独报告 \feature{mag__time__rms}。\\
Compact Health & XJTU: magnitude mean/std；PHM: \feature{h__time__mean_abs}, \feature{h__time__std}, \feature{h__time__rms} & 数据集感知通道选择，避免把 PHM horizontal 结论直接迁移到 XJTU cross-condition。\\
Compact Early Fault & \feature{mag__time__mean}, \feature{mag__time__mean_abs}, \feature{mag__time__std}, \feature{h__time__mean_abs}, \feature{v__time__std} & XJTU 加 C 级 spectral/shock 特征做 ablation；PHM 主线仍用 amplitude。\\
Reference ablation & \feature{mag__time__rms} & 只用于 sanity/reference 或单独消融，不作为独立发现。\\
\bottomrule
\end{tabularx}
\end{table}


Full feature set 是 \config{manual_basic} 全 45 个特征。

它适合作为默认输入。

Compact subset 的作用是提高解释性和减少冗余。

RUL compact subset 应优先保留 non-label-source amplitude features。

Health compact subset 应按数据集选择 channel。

Early Fault compact subset 应保留 A/B 主线，并把 C 级特征用于 ablation。

\subsection{冗余与消融}

很多 top features 是同一物理现象的重复表达。

例如：

\begin{itemize}[leftmargin=2em]
  \item \feature{mag__time__mean} 与 \feature{mag__time__mean_abs} 在当前 magnitude 构造下几乎等价；
  \item \feature{h__time__rms} 与 \feature{h__time__std} 高度相关；
  \item tsfresh 的 RMS/std/variance 与 manual features 重复；
  \item \feature{mag__time__rms} 和其他 magnitude amplitude 特征相关，但它同时是 label-source。
\end{itemize}

因此，baseline 不应只报告“全特征最好”。

更有解释价值的做法是报告：

\begin{itemize}[leftmargin=2em]
  \item full \config{manual_basic}；
  \item compact non-label-source subset；
  \item compact + Reference；
  \item dataset-specific C-level ablation。
\end{itemize}

这样才能区分“真实可泛化特征”与“标签构造参考特征”。
